% ====================================================================
% VİZE RAPOR ŞABLONU – Türk Masalları Derlemesi Web Sayfası
% --------------------------------------------------------------------
% DERLEME: Overleaf'te Menu > Compiler > XeLaTeX (pdfLaTeX de çalışır).
%
% KURALLAR (şablonu kullanmadan önce okuyun):
%  1. Bölüm başlıklarını SİLMEYİN ve sırasını DEĞİŞTİRMEYİN.
%  2. [KÖŞELİ PARANTEZ] içindeki ve gri renkli açıklama metinlerini
%     kendi bilgilerinizle değiştirin; açıklama metinlerini silin.
%  3. MEBİS'e YALNIZCA bu raporun PDF'ini yükleyin.
%     ZIP ve HTML dosyası YÜKLEMEYİN. Tüm kontroller GitHub'dan yapılır.
%  4. Kapaktaki GitHub ve YouTube bağlantıları doğru ve tıklanabilir
%     olmalıdır. Teslimden önce PDF'teki bağlantılara tıklayıp deneyin.
%  5. Haftalık katkı tablosundaki görseller için resim dosyalarını
%     (hafta1.png, hafta2.png, ...) Overleaf projesine yükleyin.
%     Dosya adlarında boşluk ve Türkçe karakter kullanmayın.
%     Resim yoksa tabloda "görsel ekleyin" kutusu görünür.
%  6. Yapay zekâ kullandıysanız Bölüm 5'i doldurun.
% ====================================================================
\documentclass[12pt,a4paper]{article}

\usepackage{iftex}
\ifPDFTeX
  \usepackage[T1]{fontenc}
  \usepackage[utf8]{inputenc}
\else
  \usepackage{fontspec}
\fi

\usepackage[turkish]{babel}
\usepackage[margin=2.4cm]{geometry}
\usepackage{setspace}
\usepackage{graphicx}
\usepackage{booktabs}
\usepackage{longtable}
\usepackage{array}
\usepackage{colortbl}
\usepackage{xcolor}
\usepackage{enumitem}
\usepackage{listings}
\usepackage{amssymb}
\usepackage{caption}
\usepackage{fancyhdr}
\usepackage{hyperref}

\onehalfspacing
\setlength{\parindent}{0pt}
\setlength{\parskip}{6pt}
\setlength{\headheight}{15pt}
\renewcommand{\arraystretch}{1.2}
\setlist{itemsep=3pt,topsep=5pt}

\definecolor{ana}{HTML}{514563}
\definecolor{acik}{HTML}{F2EAF8}
\definecolor{pembe}{HTML}{FFF2F7}
\definecolor{mavi}{HTML}{E6F2FC}

\hypersetup{
  colorlinks=true,
  linkcolor=ana,
  urlcolor=ana,
  pdftitle={Kalbur Saman - Vize Raporu},
  pdfauthor={Fatmanur Konya}
}

\lstset{
  basicstyle=\ttfamily\small,
  backgroundcolor=\color{acik},
  frame=single,
  rulecolor=\color{ana},
  breaklines=true,
  columns=fullflexible,
  keepspaces=true,
  showstringspaces=false
}

\pagestyle{fancy}
\fancyhf{}
\fancyhead[L]{Kalbur Saman}
\fancyhead[R]{Vize Raporu}
\fancyfoot[C]{\thepage}

% BAGLANTILAR HAZIR OLDUGUNDA BU IKI SATIR DEGISTIRILECEK.
% Ornek:
% \newcommand{\githublink}{\url{https://github.com/kullanici/depo}}
% \newcommand{\youtubelink}{\url{https://www.youtube.com/watch?v=KIMLIK}}

\newcommand{\githublink}{https://github.com/fatmanurkonya/PROGRAMLAMAYA-G-R-V-ZE-G-Z/tree/main}
\newcommand{\youtubelink}{\url{https://youtu.be/banDF5_eRGo}}

\newcommand{\universite}{
İstanbul Medipol Üniversitesi\\
Güzel Sanatlar Tasarım ve Mimarlık Fakültesi\\
Görsel İletişim Tasarımı
}

\newcommand{\dersadi}{Programlamaya Giriş}
\newcommand{\hocaadi}{Meral Şentürk}
\newcommand{\ogrenciadi}{Fatmanur Konya}
\newcommand{\ogrencino}{86250017}
\newcommand{\projeadi}{Kalbur Saman}

% ISARETLI KUTU
\newcommand{\tik}{$\boxtimes$}

% GORSEL KOMUTU
% Dosya yuklenmisse gosterir, yuklenmemisse yerini belirtir.
\newcommand{\raporgorsel}[3]{
  \begin{figure}[htbp]
    \centering
    \IfFileExists{#1}{
      \includegraphics[
        width=#2\linewidth,
        height=0.67\textheight,
        keepaspectratio
      ]{#1}
    }{
      \fbox{
        \parbox{0.86\linewidth}{
          \centering
          \vspace{12pt}
          Ekran görüntüsü Overleaf'e yüklenecek.\\
          \footnotesize\texttt{\detokenize{#1}}
          \vspace{12pt}
        }
      }
    }
    \caption{#3}
  \end{figure}
}

\newcommand{\haftagorsel}[1]{
  \IfFileExists{#1}{
    \includegraphics[
      width=\linewidth,
      height=3.4cm,
      keepaspectratio
    ]{#1}
  }{
    \fbox{
      \parbox{0.88\linewidth}{
        \centering\footnotesize
        Görsel yüklenecek.\\
        \texttt{\detokenize{#1}}
      }
    }
  }
}

\newcommand{\katilim}{
  \par\small
  \tik\ Ders içinde katılım sergiledim.
  \par
  $\square$ Telafi yaptım.
}

\begin{document}

% =========================================================
% KAPAK
% =========================================================

\begin{titlepage}
\centering
\vspace*{0.7cm}

{\large\universite\par}

\vspace{1.4cm}

{\Huge\bfseries\color{ana}\projeadi\par}

\vspace{0.4cm}

{\Large Türk Masalları Derlemesi Web Sayfası\par}
{\large Vize Raporu\par}

\vspace{1.3cm}

{\large\textbf{Ders:} \dersadi\par}
{\large\textbf{Öğretim Üyesi:} \hocaadi\par}

\vspace{1cm}

{\large\textbf{Hazırlayan:} \ogrenciadi\par}
{\large\textbf{Öğrenci No:} \ogrencino\par}
{\large\textbf{Sınıf:} 2. sınıf\par}

\vspace{1.2cm}

\fbox{
  \parbox{0.88\linewidth}{
    \centering
    \vspace{8pt}
    {\large\bfseries Proje Bağlantıları}\\[9pt]
    \textbf{GitHub deposu:}\\
    \githublink\\[9pt]
    \textbf{YouTube sunum videosu:}\\
    \youtubelink
    \vspace{8pt}
  }
}

\vfill
{\large 2026--2027 Güz Dönemi\par}

\end{titlepage}

\tableofcontents
\newpage

% =========================================================
\section{Giriş}
% =========================================================

Kalbur Saman, sekiz masalın kısa uyarlamalarını bir araya
getiren bir web sayfasıdır. Projenin amacı, kullanıcının
istediği masalı kolayca bulabilmesini ve sade bir okuma
alanında okuyabilmesini sağlamaktır. Sayfanın tasarımında
pembe, açık mavi ve lila tonları tercih edilmiştir.
Masal atmosferini desteklemek için karşılama bölümünde
bir şato görseline yer verilmiştir. İçerik, tasarım ve
etkileşimler tek HTML dosyasında birleştirilmiştir.
Bu raporda masal seçimi, tasarım kararları, teknik yapı,
yapay zekâ kullanımı ve beş haftalık ders içi çalışmalar
açıklanmaktadır.

\subsection{Proje Özeti}

Sayfanın üst bölümünde proje adı, gezinme bağlantıları,
karşılama metni ve kapak görseli bulunmaktadır.
``Masal sandığı'' bölümünde sekiz masal kartı yer almaktadır.
Kullanıcı masal adına göre arama yapabilir veya kategori
seçerek listeyi filtreleyebilir. ``Masal oku'' düğmesi
seçilen masalı bir okuma penceresinde açmaktadır.
Bu pencerede yazı boyutu değiştirilebilir ve pencere
kapatılarak masal listesine dönülebilir.

\subsection{Gereksinimlerin Karşılanması}

Aşağıdaki kutular teslim hedeflerini gösterecek şekilde
işaretlenmiştir. Son cihaz testleri, metin düzenlemeleri
ve bağlantı kontrolleri teslim öncesinde tamamlanacaktır.

\begin{center}
\begin{tabular}{@{}p{10.7cm}c@{}}
\toprule
\textbf{Gereksinim} & \textbf{Durum}\\
\midrule
Masal sayısı 8--12 arasında: 8 masal & \tik\\
Tek HTML dosyası & \tik\\
Stil kodları içsel: \texttt{<style>} içinde & \tik\\
Betik kodları içsel: \texttt{<script>} içinde & \tik\\
İnternetsiz açılıp çalışma & \tik\\
JavaScript'in anlamlı işlevlerde kullanılması & \tik\\
Masal metinlerinin kendi cümlelerimle düzenlenmesi & \tik\\
Sayfanın telefon ekranına uyum sağlaması & \tik\\
\bottomrule
\end{tabular}
\end{center}

Mevcut masal taslakları yapay zekâ desteğiyle hazırlanmıştır.
Kendi cümlelerimle düzenleme maddesi, teslim öncesinde
yapılacak kişisel metin düzenlemesini de kapsamaktadır.
Görselin dosyaya gömüldüğü sürümün internetsiz çalışması
cihaz üzerinde ayrıca kontrol edilecektir.

% =========================================================
\section{Masalların Seçimi}
% =========================================================

Masal listesi; Keloğlan, sihirli nesneler ve hayvan
karakterleri çevresinde çeşitlilik oluşturacak şekilde
hazırlanmıştır. Sayfadaki tür adları, kullanıcıların
filtreleme yapabilmesi için kullanılan kategorilerdir.

\begin{longtable}{@{}cp{3.9cm}p{2.8cm}p{5.1cm}@{}}
\toprule
\textbf{No} & \textbf{Masal} & \textbf{Tür}
& \textbf{Seçilme nedeni}\\
\midrule
\endhead

1 & Keloğlan ve Sihirli Taş & Keloğlan
& İyilik ve paylaşma temasını işlemesi.\\

2 & Keloğlan ve Dev & Keloğlan
& Zekâ ile güç arasındaki karşılaşmayı anlatması.\\

3 & Sabır Taşı & Sihirli masallar
& Sabır ve cesaret temalarına yer vermesi.\\

4 & Altın Bülbül & Sihirli masallar
& Yolculuk ve arayış duygusu oluşturması.\\

5 & Sihirli Sofra & Sihirli masallar
& Paylaşma ve bereket temasını desteklemesi.\\

6 & İnatçı Keçi & Hayvan masalları
& Uzlaşma ve inatçılık üzerine düşünme sağlaması.\\

7 & Kurnaz Tilki & Hayvan masalları
& Aldatıcı sözlere karşı dikkatli olmayı işlemesi.\\

8 & Açıl Sofram Açıl & Sihirli masallar
& Sihirli bir nesne çevresinde olay örgüsü kurulması.\\

\bottomrule
\end{longtable}

Sitedeki metinler belirli bir kitaptan birebir
aktarılmamıştır; yapay zekâ desteğiyle oluşturulan kısa
uyarlamalardır. Teslim öncesinde anlatıların kaynaklarla
karşılaştırılması ve Türk masalları kapsamına
uygunluğunun doğrulanması gerekmektedir. Özellikle hayvan
anlatılarının kökenleri doğrulanmadan yalnızca
Türk geleneğine ait oldukları ileri sürülmemektedir.

% =========================================================
\section{Tasarım Kararları}
% =========================================================

\subsection{Konuya Uygunluk}

``Kalbur Saman'' adı, masal tekerlemelerini çağrıştırdığı
için seçilmiştir. ``Bir varmış, bir yokmuş'' başlığı,
sayfanın konusunu ilk bakışta anlatmaktadır.
Kapaktaki şato görseli fantastik bir atmosfer
oluşturmak amacıyla kullanılmıştır. Görsel, Türk
masallarına ilişkin tarihî bir belge olarak
sunulmamaktadır.

Açık renkli zemin ve yuvarlatılmış kartlar,
sayfanın yumuşak bir görünüme sahip olmasını
sağlamaktadır. Masal kartlarında aynı düzenin
kullanılmasıyla içerikler arasında görsel
bütünlük oluşturulmuştur.

\subsection{Renk Paleti ve Yazı Tipi}

\begin{center}
\begin{tabular}{@{}p{2.6cm}p{2.8cm}p{7.1cm}@{}}
\toprule
\textbf{Ad} & \textbf{Kod / tür}
& \textbf{Kullanıldığı alan ve gerekçe}\\
\midrule

Açık pembe & \texttt{\#FFF2F7}
& Sayfanın ana zemininde kullanılmıştır.\\

Koyu mor & \texttt{\#514563}
& Açık zemindeki metinlerin belirginleşmesini sağlar.\\

Açık mavi & Pastel ton
& Görsel çerçevesi ve bilgi alanlarında kullanılır.\\

Lila & Pastel ton
& Düğme ve kart ayrıntılarını vurgular.\\

Georgia & Serif
& Başlıklarda ve masal metinlerinde kitap hissi oluşturur.\\

Arial & Sans serif
& Menü ve yardımcı arayüz metinlerinde kullanılır.\\

\bottomrule
\end{tabular}
\end{center}

Kahverengi yerine pembe, açık mavi ve lila tonlarına
ağırlık verilmiştir. Açık zemin üzerinde koyu renkli
metin kullanılarak okunabilirlik desteklenmiştir.

\subsection{Okunabilirlik, Gezinme ve Duyarlı Tasarım}

\begin{itemize}

\item \textbf{Okunabilirlik:}
Masal metinleri paragraflara ayrılmıştır.
Okuma penceresindeki A-- ve A+ düğmeleri,
yazı boyutunu değiştirmektedir.

\item \textbf{Gezinme:}
Üst menü, arama kutusu ve kategori seçimi,
istenen masala ulaşmayı kolaylaştırmaktadır.

\item \textbf{Duyarlı tasarım:}
Bilgisayarda iki sütunda gösterilen kartlar,
dar ekranlarda tek sütuna dönüşmektedir.

\item \textbf{Tutarlılık:}
Kartların başlık, açıklama ve düğme yerleşimleri
aynı düzeni takip etmektedir.

\end{itemize}

\subsection{Ekran Görüntüleri}

Görseller, sayfanın bilgisayar ve telefon üzerindeki
görünümünü belgelemektedir.

\raporgorsel
{Ekran Resmi 2026-11-04 19.43.00.png}
{0.98}
{Bilgisayarda ana sayfa, karşılama başlığı ve kapak görseli.}

\raporgorsel
{Ekran Resmi 2026-11-04 19.43.07.png}
{0.98}
{Arama alanı, kategori menüsü ve ilk masal kartları.}

\raporgorsel
{Ekran Resmi 2026-11-04 19.43.48.png}
{0.98}
{Diğer masal kartları, hakkında alanı ve alt bilgi.}

\raporgorsel
{Ekran Resmi 2026-11-04 19.43.59.png}
{0.98}
{Masal okuma penceresi ve yazı boyutu düğmeleri.}

\raporgorsel
{IMG_0490.PNG}
{0.42}
{Telefon ekranında masal kartlarının tek sütunlu görünümü.}

\raporgorsel
{IMG_0494.PNG}
{0.42}
{Telefon ekranında hakkında alanı ve alt bilgi.}

\clearpage

% =========================================================
\section{Gerçekleştirme}
% =========================================================

\subsection{HTML Yapısı}

Sayfanın üst bölümünde proje adı ile ``Masallar''
ve ``Hakkında'' bağlantıları bulunmaktadır.
Ana içerikte karşılama alanı, kapak görseli,
arama kutusu, kategori seçimi ve masal kartları
yer almaktadır. Alt bölümde proje açıklaması,
görsel kaynağı ve öğrenci bilgileri bulunmaktadır.

Masalların okunması için \texttt{dialog} etiketi
kullanılmıştır. Seçilen masalın başlığı ve
paragrafları bu pencereye aktarılmaktadır.
Böylece her masal için ayrı bir HTML dosyası
oluşturmaya gerek kalmamıştır.

\subsection{İçsel CSS}

CSS kuralları HTML dosyasının \texttt{<style>}
alanında bulunmaktadır. Renkler, yazı tipleri,
kart düzeni, kenarlıklar ve boşluklar bu alanda
tanımlanmıştır. Kartların yerleşiminde grid,
dar ekranlara uyum sağlamak için medya
sorgusu kullanılmıştır.

Aşağıdaki örnek, projede kullanılan iki sütundan
tek sütuna geçiş yaklaşımını sadeleştirerek
göstermektedir. Sınıf adı örnek amaçlıdır.

\begin{lstlisting}[caption={Duyarlı kart düzeninin sadeleştirilmiş örneği}]
.masallar {
  display: grid;
  grid-template-columns: repeat(2, 1fr);
  gap: 20px;
}

@media (max-width: 600px) {
  .masallar {
    grid-template-columns: 1fr;
  }
}
\end{lstlisting}

Kapak görselinin gömüldüğü sürümde görsel verisi
doğrudan \texttt{img} etiketinin \texttt{src}
alanında tutulmaktadır. Bu yöntem dosya boyutunu
artırsa da görselin internetten indirilmesi
gereğini ortadan kaldırmaktadır.

\subsection{İçsel JavaScript}

Masallar bir dizi içinde tutulmaktadır.
Her masal için ad, tür, kısa açıklama ve
metin alanları bulunmaktadır.

\texttt{listele} fonksiyonu, arama ve kategori
seçimine göre kartları oluşturmaktadır.
\texttt{masalAc} fonksiyonu, seçilen masalın
metnini okuma penceresinde göstermektedir.
Yazı boyutu düğmeleri, metin boyutunu belirlenen
alt ve üst sınırlar içinde değiştirmektedir.

Aramada Türkçe büyük ve küçük harf dönüşümü
dikkate alınmıştır. Aşağıdaki örnek, arama
mantığını sadeleştirerek göstermektedir.

\begin{lstlisting}[caption={Masal adına göre arama mantığı}]
const aranan = arama.value
  .toLocaleLowerCase("tr");

const sonuc = masallar.filter(masal =>
  masal.ad.toLocaleLowerCase("tr")
    .includes(aranan)
);
\end{lstlisting}

\subsection{Sayfanın Özellikleri}

\begin{itemize}
\item Sekiz masalı kartlar hâlinde listeleme.
\item Masal adına göre arama.
\item Kategoriye göre filtreleme.
\item Listelenen masal sayısını gösterme.
\item Seçilen masalı okuma penceresinde açma.
\item Yazı boyutunu artırma ve azaltma.
\item Okuma penceresini kapatarak listeye dönme.
\item Dar ekranlarda tek sütunlu kart düzeni.
\end{itemize}

% =========================================================
\section{Yapay Zekâ Kullanımı}
% =========================================================

Projenin hazırlanmasında ChatGPT'den kod,
metin taslağı ve düzenleme desteği alınmıştır.
Sayfadaki masal metinleri de bu destek kapsamında
oluşturulan kısa uyarlamalardır.

\begin{longtable}{@{}p{1.7cm}p{3.2cm}p{4.6cm}p{3.5cm}@{}}
\toprule
\textbf{Araç} & \textbf{Kullanım amacı}
& \textbf{İstem özeti} & \textbf{Kullanıldığı yer}\\
\midrule
\endhead

ChatGPT
& HTML, CSS ve JavaScript desteği
& Sekiz masallı, pastel renkli; arama, filtreleme
ve okuma penceresi bulunan tek dosya hazırlanması.
& Sayfanın yapısı ve etkileşimleri.\\

ChatGPT
& Renk ve kod düzenleme
& Kahverenginin kaldırılması; pembe ve açık mavi
tonlarının kullanılması, kodun kısaltılması.
& CSS ve genel düzen.\\

ChatGPT
& Masal metni taslakları
& Belirlenen başlıklara uygun kısa uyarlamalar
hazırlanması.
& Masal okuma penceresi.\\

ChatGPT
& Çevrim dışı kullanım düzenlemesi
& Kapak görselinin HTML dosyasına gömülmesi.
& Görselin \texttt{src} alanı.\\

ChatGPT
& Rapor taslağı
& Verilen bilgiler ve ekran görüntüleriyle
öğretim üyesinin şablonunun doldurulması.
& Rapor metni ve LaTeX düzeni.\\

\bottomrule
\end{longtable}

\subsection{Anladığım ve Kendim Yaptığım Kısımlar}

Proje adını, renk tercihlerini ve sayfanın genel
görünümünü belirledim. Görselin büyüklüğü ve
renkler konusunda düzenleme taleplerinde bulundum.
Hazırlanan sayfayı HTML Online Viewer üzerinden
görüntüledim ve değiştirilmesini istediğim
noktaları belirttim. Bilgisayar ve telefon
görünümlerinin ekran görüntülerini hazırladım.

Kod hazırlanırken yapay zekâ desteği alındığı için
bütün kodu tek başıma yazdığım iddiasında bulunmuyorum.
HTML'nin içeriği, CSS'nin görünümü ve JavaScript'in
etkileşimleri yönettiğini çalışma üzerinden
inceledim. Sunum öncesinde arama, filtreleme,
masal açma ve yazı boyutu değiştirme işlevlerinin
mantığını yeniden gözden geçireceğim.

% =========================================================
\section{Haftalık Katkılar}
% =========================================================

Bu bölümde beş haftalık ders içi çalışmalar
sunulmuştur. Çalışmalar telafi olarak yapılmamıştır.
Hafta sıralaması 23 Eylül, 30 Eylül, 7 Ekim,
14 Ekim ve 21 Ekim şeklindedir.

GitHub deposu henüz hazırlanmadığından geçmiş
çalışmalar için doğrulanmamış commit bilgisi
yazılmamıştır. Çalışmalar mevcut ekran
görüntüleriyle belgelenmiştir.

\begin{longtable}{@{}
>{\raggedright\arraybackslash}p{1.5cm}
>{\raggedright\arraybackslash}p{6.3cm}
>{\centering\arraybackslash}p{5.7cm}@{}}
\toprule
\textbf{Hafta} & \textbf{Yapılan çalışma}
& \textbf{Görsel}\\
\midrule
\endhead

\textbf{1. hafta}\newline
23 Eylül
&
WorldWideWeb başlıklı sayfa üzerinde temel HTML
yapısı çalışıldı. Başlık, paragraf, liste ve
bağlantılar kullanılarak içerik bölümlere ayrıldı.
Sayfa içi bağlantılarla ilgili bölüme ulaşma
yaklaşımı incelendi.

\textbf{Commit:} Henüz eklenmedi.
\katilim
&
\haftagorsel{resim-20261007-095109-1a9ba975_9055.jpeg}
\\
\midrule

\textbf{2. hafta}\newline
30 Eylül
&
``Kimler Geldi Kimler Geçti'' dizisi için tanıtım
sayfası hazırlandı. Konu, özet, oyuncular,
görseller ve açılır soru alanları üzerinde çalışıldı.
Ayrıca tıklama, çift tıklama, fareyle üzerine
gelme ve bağlantı açma örnekleri yapıldı.

\textbf{Commit:} Henüz eklenmedi.
\katilim
&
c
\haftagorsel{resim-20260927-190456-d1e387c4_5722.jpeg}
\haftagorsel{resim-20260927-190536-69940d42_3828.jpeg}
\par\vspace{6pt}
\\
\midrule

\textbf{3. hafta}\newline
7 Ekim
&
Kullanıcıdan isim alan bir giriş kutusu ve düğme
üzerinde çalışıldı. JavaScript ile giriş alanının
değeri alınarak sayfada selamlama metni oluşturma
yaklaşımı incelendi. Boş giriş için koşul
kullanımı çalışıldı.

\textbf{Commit:} Henüz eklenmedi.
\katilim
&
\haftagorsel{Ekran Resmi 2026-09-30 13.39.24 (1).png}
\\
\midrule
\haftagorsel{Ekran Resmi 2026-09-30 14.55.14 (1).png}
\par\vspace{10pt}
\\
\midrule



\textbf{4. hafta}\newline
14 Ekim
&
Bağlantı ve düğme ile farklı sayfalara yönlendirme
çalışıldı. Renk eşleştirme örneğinde kartların
veri özellikleri okundu; aynı ve farklı renklerin
karşılaştırılması incelendi. Zamanlayıcı ve
sınıf ekleme işlemleri kullanıldı.

\textbf{Commit:} Henüz eklenmedi.
\katilim
&
\haftagorsel{Ekran Resmi 2026-10-14 14.02.55.png}
\par\vspace{5pt}
\haftagorsel{Ekran Resmi 2026-09-30 14.30.04 (1).png}
\\
\midrule

\textbf{5. hafta}\newline
21 Ekim
&
Harici CSS dosyası ile HTML arasındaki bağlantı
üzerinde çalışıldı. Kiehl's temalı sayfada
başlık, gezinme, bölüm ve kart düzenleri ele alındı.
Bootstrap sınıflarıyla düzenleme örnekleri yapıldı.
Düğmeyle bir kutunun rengini değiştirme
çalışması gerçekleştirildi.

\textbf{Commit:} Henüz eklenmedi.
\katilim
&
\haftagorsel{Ekran Resmi 2026-10-21 13.40.26.png}
\par\vspace{5pt}
\haftagorsel{Ekran Resmi 2026-11-04 20.25.31.png}
\\
\bottomrule
\end{longtable}

Harici CSS ve Bootstrap örnekleri ders içi
alıştırmalardır. Vize projesinde tek HTML dosyası
koşuluna uygun olarak içsel CSS ve içsel
JavaScript kullanılmıştır.

% =========================================================
\section{Test ve Değerlendirme}
% =========================================================

Bilgisayar ekran görüntüleri Safari üzerinden
alınmıştır. Telefon ekran görüntülerinde kartların
tek sütuna dönüştüğü ve alt bölümlerin görüntülendiği
görülmektedir. Görünüm kontrolleri ile etkileşim
testleri ayrı değerlendirilmiştir.

\begin{longtable}{@{}p{3.7cm}p{4.6cm}p{4.7cm}@{}}
\toprule
\textbf{Özellik} & \textbf{Bilgisayar}
& \textbf{Telefon}\\
\midrule
\endhead

Sayfa yerleşimi
& Ekran görüntüsü mevcut.
& Tek sütun görünümü mevcut.\\

Masal kartları
& Kartlar görüntüleniyor.
& Kartlar görüntüleniyor.\\

Masal penceresi
& Açık pencere görüntüsü mevcut.
& Son işlev kontrolü yapılacak.\\

Arama ve kategori
& Kod düzeyindeki kontroller geçti.
& Cihaz üzerinde son kontrol yapılacak.\\

Yazı boyutu
& Kod düzeyindeki sınır kontrolleri geçti.
& Cihaz üzerinde son kontrol yapılacak.\\

İnternetsiz kullanım
& Gömülü görselli sürüm hazır; bağlantı kapalı
son kontrol yapılacak.
& Bağlantı kapalı son kontrol yapılacak.\\

\bottomrule
\end{longtable}

\textbf{Görüntü alınan tarayıcı:} Bilgisayarda Safari.

Kod düzeyindeki kontrollerde sekiz masalın
listelenmesi, metinlerin açılması, Türkçe arama,
kategori seçimi ve yazı boyutu sınırları
incelenmiştir. Bu kontroller gerçek cihazdaki
son kullanım testinin yerine geçmemektedir.

% =========================================================
\section{Sunum Videosu}
% =========================================================

Sunum videosu hazırlanarak bağlantısı bu bölüme
ve kapaktaki bağlantı alanına eklenecektir.

\begin{center}
\begin{tabular}{@{}p{3.7cm}p{8.6cm}@{}}
\textbf{Video süresi:}
& Kayıt sonrasında eklenecek; 12--18 dakika.\\
\textbf{Görünürlük:}
& Listelenmemiş veya herkese açık olarak ayarlanacak.\\
\textbf{Bağlantı:}
& \youtubelink\\
\end{tabular}
\end{center}

Sunumda projenin amacı, sayfanın kullanımı,
HTML yapısı, CSS kararları ve JavaScript
işlevleri açıklanacaktır. Ders içi çalışmalar
ve yapay zekâ desteği de belirtilecektir.

\subsection{Video Kontrol Listesi}

Aşağıdaki kutular teslim öncesi tamamlanacak
kontrol hedefleri olarak işaretlenmiştir.

\begin{itemize}[label={}]
\item \tik\ Süre 12--18 dakika aralığında olacak.
\item \tik\ Sunum boyunca kamera açık ve yüzüm görünür olacak.
\item \tik\ Video listelenmemiş veya herkese açık olacak.
\item \tik\ Bağlantı gizli pencerede açılıp denenecek.
\item \tik\ YouTube'un video işlemeyi tamamladığı kontrol edilecek.
\item \tik\ Ses ve kamera baştan sona kontrol edilecek.
\item \tik\ Kapaktaki bağlantı ile bu bölümdeki bağlantı aynı olacak.
\end{itemize}

% =========================================================
\section{Sonuç ve Kendi Değerlendirmem}
% =========================================================

Kalbur Saman projesinde sekiz kısa masal
uyarlamasını bir araya getiren bir web sayfası
hazırlandı. HTML ile içerik, CSS ile görünüm
ve JavaScript ile kullanıcı etkileşimleri
birleştirildi. Pastel renkler ve kart düzeniyle
bütünlüklü bir görünüm oluşturulmaya çalışıldı.
Çalışma sırasında özellikle görsel boyutu,
renklerin düzenlenmesi ve kodun sadeleştirilmesi
üzerinde duruldu.

Arama, kategori seçimi ve okuma penceresi,
sayfanın etkileşimli hâle gelmesini sağladı.
Telefon görünümü için kart düzeni tek sütuna
uyarlandı. Yapay zekâ desteği raporda açıklandı.
Teslim öncesinde kaynak kontrolü, kişisel metin
düzenlemesi, internetsiz cihaz testleri,
GitHub deposu ve sunum videosu tamamlanacaktır.

% =========================================================
\section{Kaynakça}
% =========================================================

\begin{enumerate}

\item
Programlamaya Giriş dersi.
\textit{Türk Masalları Derlemesi Web Sayfası Vize Briefi}.
Öğretim üyesi tarafından paylaşılan PDF.

\item
Programlamaya Giriş dersi.
\textit{Vize Rapor Şablonu}.
Öğretim üyesi tarafından paylaşılan
\texttt{Vize\_Sablon} arşivindeki LaTeX dosyası.

\item
Stöckel, Andrea.
\textit{Fantasy Fairytale Castle}.
Public Domain Pictures.
Sayfada kullanılan kapak görseli.
\url{https://www.publicdomainpictures.net/en/view-image.php?image=578308&picture=fantasy-fairytale-castle}

\item
OpenAI.
\textit{ChatGPT}.
Kod, kısa masal uyarlamaları ve rapor taslağı
hazırlanmasında kullanılan yapay zekâ aracı.
\url{https://chatgpt.com/}

\end{enumerate}



% =========================================================
\appendix
\section{Teslim Kontrol Listesi}
% =========================================================

Kutular, teslim öncesinde tamamlanacak işlemleri
gösterecek şekilde işaretlenmiştir.

\begin{itemize}[label={}]

\item \tik\
MEBİS'e yalnızca rapor PDF'i yüklenecek;
ZIP ve HTML dosyası yüklenmeyecek.

\item \tik\
PDF dosya adı
\texttt{86250017\_FatmanurKonya\_vize.pdf}
olarak düzenlenecek.

\item \tik\
GitHub ve YouTube bağlantıları eklenecek,
PDF üzerinden tıklanarak denenecek.

\item \tik\
GitHub deposuna tek HTML dosyası, rapor PDF'i,
LaTeX dosyası ve rapor görselleri eklenecek.

\item \tik\
README dosyasına proje açıklaması,
kullanım bilgileri ve video bağlantısı eklenecek.

\item \tik\
GitHub deposu öğretim üyesinin erişimine açılacak.

\item \tik\
Son dosyalar teslim saatinden önce depoya gönderilecek.

\item \tik\
Sayfada içsel CSS ve içsel JavaScript kullanılacak.

\item \tik\
Hazır site şablonu kullanılmayacak;
yapay zekâ desteği açıklanacak.

\item \tik\
Masal kaynakları ve metinlerin ödev kapsamına
uygunluğu doğrulanacak.

\item \tik\
İnternet kapalıyken görsel, arama,
filtreleme ve okuma işlevleri denenecek.

\item \tik\
Çalışmanın mantığı sunumda açıklanacak.

\end{itemize}

\end{document}
